\documentclass{article}
\usepackage{lineno,hyperref}
\usepackage[T1]{fontenc}
%\usepackage{cite}
\usepackage{graphicx,amsmath}
\usepackage{soul,color}
\usepackage{siunitx}
\usepackage{mhchem}
\usepackage[sorting=none,citestyle=numeric-comp]{biblatex}
\bibliography{refs}
\usepackage{lineno}
%\linenumbers
\usepackage[bottom=2cm, top=2cm, right=2cm, left=2.5cm]{geometry}
\title{hexatic graphene}
\author{Me}
\date{\today}

\begin{document}

\maketitle

\section{Correlation function}

The Hamiltonian of the system is given by~\cite{Yazyev2010}:

$$H = \frac{\epsilon}{\rho}\sum_i(\theta_i^2) + \frac{\gamma}{3\sqrt{\rho}}\sum_i \sum_{j\in N_i} (|\theta_i - \theta_j|)(A-\ln{(|\theta_i - \theta_j|))}$$

where $\epsilon$ is a material-dependent constant, $\rho$ is the density of grains, $\gamma$ and $A$ are some constants. $N_i$ is the set of neighbors of grain $i$.

The partition function is given by:

$$\mathcal{Z} = \int \prod_i d\theta_i \exp{\left(-\beta \sum_i E_i\right)} = \int \prod_i d\theta_i \text{Tr}\left\{\exp{\left(-\beta H \right)}\right\}$$
where $\beta = \frac{1}{k_B T}$, $k_B$ is the Boltzmann constant and $T$ is the temperature.

\subsection{Large grain limit}

In the limit of large grains, global alignment with substrate dominates grains boundaries energy, and the misorientation distribution becomes distance-independent

$$\mathcal{Z} = \int \prod_i d\theta_i \exp{\left(-\beta \frac{\epsilon}{\rho}\sum_i \theta_i^2\right)} = \prod_i \int d\theta_i \exp{\left(-\beta \frac{\epsilon}{\rho}\theta_i^2\right)}$$

So we have a gaussian distribution for each grain orientation $\theta_i$ with mean $0$ and variance $\sigma^2 = \frac{\rho}{2\beta\epsilon}$. 

The correlation function is given by:

$$C(\theta_i, \theta_j) = \langle \phi_i \phi_j^\star \rangle = \langle e^{i6(\theta_i - \theta_j)} \rangle$$

As the angles are independent, we have:

$$C(\theta_i, \theta_j) = \left|\langle e^{i6\theta_i} \rangle \right|^2 = \left| e^{-18\sigma^2} \right|^2 = e^{-18\frac{\rho}{\beta\epsilon}}$$


\subsection{Small grain limit}

In the limit of small grains, grain boundaries energy dominates global alignment with substrate. To compute the correlation function, the Kosterlitz-Thouless and Wegner~\cite{Wegner1967} approaches are used.

The goal is to approximate the hamiltonian to a quadratic form. However, the function 
$$|\theta_i - \theta_j|(A-\ln{(|\theta_i - \theta_j|)})$$ 
is not differentiable at $0$. To overcome this issue, we approximate it by a differentiable function $f(x)$ such that $f(x) \approx x(A-\ln{(x)})$ for $x \gg 0$ and $f(x) \approx 0$ for $x \approx 0$. A possible choice is 
$$f(x) = \sqrt{x^2 + a^2}\left(A-\ln{\left(\sqrt{x^2 + a^2}\right)}\right)$$ 
where $a$ is a small parameter.

We have:

$$H = \frac{\epsilon}{\rho}\sum_i\theta_i^2 + \frac{\gamma}{3\sqrt{\rho}}\sum_i \sum_{j\in N_i} \left(\sqrt{|\theta_i - \theta_j|^2 + a^2}\right)\left(A-\frac{1}{2}\ln{\left(|\theta_i - \theta_j|^2 + a^2\right)}\right)$$

Let $g_a (\Delta \theta)$ be the function defined by:

$$g_a (\Delta \theta) = \left(\sqrt{|\Delta \theta|^2 + a^2}\right)\left(A-\frac{1}{2}\ln{\left(|\Delta \theta|^2 + a^2\right)}\right)$$

For a small $\Delta \theta$ and assume $\Delta \theta \geq 0$, we can expand $g_a$:

$$g_a (\Delta \theta) = cste + \frac{\Delta \theta^2}{2a}\left(A - \ln{a} - 1\right) + \mathcal{O}\left(\Delta \theta^4\right)$$

So the grain boundary term can be approximated by a quadratic term for small misorientation angles. The hamiltonian becomes:

$$H = \int d^2r \left[\frac{\mu}{2} \theta^2 + \frac{K}{2}(\nabla \theta)^2\right]$$
where $\mu = \frac{2\epsilon}{\ell^2 \rho}$, $K = \frac{\gamma}{3\sqrt{\rho}}\frac{z}{2a}\left(A - \ln{a} - 1\right)$, $\ell$ is the average grain size and $z$ is the coordination number of the lattice.

This is exactly an Hamiltonian of a massive gaussian field. In this case, we have:

$$\left\langle \left(\theta(\textbf{r}) - \theta(0)\right)^2 \right\rangle = \frac{k_B T}{2 \pi K} \left(\ln{\frac{r}{\xi}} + K_0 \frac{r}{\xi}\right)$$
where $\xi = \sqrt{\frac{K}{\mu}}$ is the correlation length.

The correlation function is given by:

$$G_6(r) = \left\langle e^{i6(\theta(\textbf{r}) - \theta(0))} \right\rangle = e^{-18\left\langle \left(\theta(\textbf{r}) - \theta(0)\right)^2 \right\rangle}$$

So we have:

$$G_6(r) = 
\begin{cases} 
      \left(\frac{r}{\ell}\right)^{-\eta_6} & \ell \ll r \ll \xi \\
      \left(\frac{\xi}{\ell}\right)^{-\eta_6} e^{-\eta_6 K_0 r/\xi} & r \gg \xi \\
\end{cases}
$$
where $\eta_6 = \frac{18 k_B T}{2 \pi K}$.

\printbibliography
\end{document}